% !TEX program = xelatex
% !TEX root = ../../TFM.tex
\documentclass[../../TFM.tex]{subfiles}

\begin{document}


	% =====================================================
	% RESULTADOS / HALLAZGOS / DESARROLLO
	% =====================================================
	\chapter{Resultados}
	\label{chap:resultados}
	\thesisstartchap{resultados}

	\section{Muestra analítica y especificaciones estimadas}
	\label{sec:muestra-analitica-resultados}

	Esta sección presentará de manera sintética la muestra efectivamente utilizada en las estimaciones econométricas, una vez completada la construcción definitiva del panel. A diferencia del capítulo de datos, donde se describe la base y se documentan sus principales patrones descriptivos, aquí solo se reportarán los elementos necesarios para interpretar los resultados: número de países y observaciones, cobertura temporal efectiva, distribución por tipo de recurso predominante y cambios derivados de la disponibilidad de información para las variables del modelo.

	La muestra corresponderá al umbral único de dependencia externa de recursos $\theta=20\%$, que identifica 55 países antes de cruzar la lista con las demás variables econométricas. El período econométrico será 1996--2021. Se reportará con claridad cualquier reducción del número de países u observaciones causada por la cobertura efectiva de las variables, sin redefinir la regla de selección ni imputar datos faltantes.


	\section{Resultados principales: complejidad económica}
	\label{sec:resultados-eci}

	Esta sección presentará los resultados principales de la tesis, utilizando el Índice de Complejidad Económica ($ECI$) como variable dependiente. En estas estimaciones se evaluará la asociación entre la intensidad macroeconómica de las rentas extractivas ($RENTS$), la calidad institucional ($INST$), su interacción y los canales productivos, estructurales, macroeconómicos, financieros y de capacidades definidos en el capítulo metodológico.

	La interpretación se organizará en torno a tres elementos: el signo y magnitud de los coeficientes estimados, la relevancia sustantiva de las asociaciones encontradas y su consistencia con las hipótesis de trabajo. En esta especificación, el índice de concentración exportadora ($HHI$) se mantendrá como regresor del canal estructural, dado que la variable dependiente es $ECI$ y no una transformación directa de dicho índice.


	\section{Resultados complementarios: diversificación exportadora}
	\label{sec:resultados-divx}

	Como ejercicio complementario, esta sección presentará estimaciones que utilizan la diversificación exportadora ($DIVX=1-HHI$) como variable dependiente. El objetivo no es reemplazar la medición principal basada en complejidad económica, sino evaluar si los mecanismos asociados a mayores niveles de sofisticación exportadora también se relacionan con una canasta exportadora menos concentrada.

	En esta especificación, $HHI$ se excluirá del conjunto de regresores, dado que $DIVX$ se construye directamente como su complemento. La lectura de estos resultados será, por tanto, complementaria: mientras el modelo con $ECI$ aproxima la sofisticación y las capacidades productivas reveladas, el modelo con $DIVX$ permite observar de manera más directa la dimensión de diversificación de la canasta exportadora.


	\section{Discusión de resultados}
	\label{sec:discusion-resultados}

	La discusión integrará la evidencia econométrica con el marco teórico desarrollado en los capítulos anteriores. En lugar de limitarse a resumir coeficientes, esta sección interpretará los resultados a la luz de los mecanismos de enfermedad holandesa, acumulación de capacidades, diversificación exportadora y restricciones institucionales que condicionan la transformación estructural externa en economías dependientes de recursos naturales.

	Adicionalmente, se contrastarán los resultados obtenidos con la literatura empírica revisada, identificando coincidencias, diferencias y posibles aportes del estudio. Esta lectura comparada permitirá valorar si la relación entre dependencia extractiva y complejidad económica es sistemática o si depende de mediaciones institucionales y productivas específicas.



	\subsection{Ranking descriptivo de complejidad y diversificación}
	\label{subsec:ranking}

	Con el fin de complementar el análisis econométrico, se propone un ejercicio descriptivo orientado a comparar la posición relativa de las economías dependientes de recursos naturales en términos de complejidad económica y diversificación exportadora. A diferencia de enfoques que miden la transformación estructural a partir del crecimiento del sector no extractivo, este trabajo utiliza el \textit{Economic Complexity Index} ($ECI$) como medida principal de capacidades productivas reveladas y $DIVX=1-HHI$ como indicador complementario de dispersión de la canasta exportadora.

	La lógica de este enfoque se basa en la literatura de complejidad económica, según la cual la diversificación no debe entenderse únicamente como la expansión de sectores no extractivos, sino como la acumulación de capacidades que permiten a una economía producir y exportar bienes más diversos y sofisticados. En este sentido, el $ECI$ captura la sofisticación de la estructura exportadora, mientras que $DIVX$ resume de forma más directa el grado de concentración o dispersión de dicha canasta.

	En consecuencia, el ejercicio propuesto considera tres lecturas complementarias: (i) el nivel de complejidad económica, que refleja el estado actual de las capacidades productivas reveladas; (ii) la dinámica de la complejidad, que captura la evolución de dichas capacidades en el tiempo; y (iii) el nivel de diversificación exportadora, que permite identificar economías cuya canasta está menos concentrada en pocos productos.

	Formalmente, se definen:

	\begin{equation}
		RTC_i^{nivel} = ECI_i
	\end{equation}

	\begin{equation}
		RTC_i^{din} = \Delta ECI_i = ECI_{i,t} - ECI_{i,t-k}
	\end{equation}

	\begin{equation}
		RTD_i = DIVX_i = 1-HHI_i
	\end{equation}

	donde $RTC_i^{nivel}$ representa la posición relativa en términos de complejidad económica, $RTC_i^{din}$ captura el cambio en dicha posición a lo largo del tiempo y $RTD_i$ resume la diversificación exportadora mediante el complemento del índice de Herfindahl--Hirschman.
	
	Esta distinción permite diferenciar entre economías que ya poseen estructuras exportadoras complejas, aquellas que muestran mejoras dinámicas en complejidad y aquellas que presentan canastas exportadoras relativamente menos concentradas. En particular, el análisis conjunto de estas dimensiones permite identificar trayectorias diferenciadas dentro del grupo de economías dependientes de recursos naturales.
	
	De este modo, el ranking no solo ordena a los países según su nivel de complejidad económica, sino que también permite clasificar su desempeño en términos de transición estructural externa, distinguiendo entre economías consolidadas, economías en proceso de diversificación y economías con persistente concentración extractiva.
	
	Finalmente, este enfoque permite identificar heterogeneidad significativa entre economías con similares niveles de dependencia de recursos, evidenciando que dicha dependencia no determina de manera unívoca la capacidad de avanzar hacia estructuras exportadoras más complejas y diversificadas. Por el contrario, algunos países logran desarrollar canastas más sofisticadas o avanzar dinámicamente hacia ellas, lo que sugiere la existencia de factores adicionales ---institucionales, educativos y tecnológicos--- que condicionan las trayectorias de transformación estructural.
	
	\thesisendchap{resultados}
	
\end{document}
